\documentclass{article}
\usepackage{amsmath,amssymb,amsthm}
\usepackage{algorithm}
\usepackage{algpseudocode}
\usepackage{bm}
\usepackage{hyperref}
\newtheorem{theorem}{Theorem}
\newtheorem{lemma}{Lemma}
\newtheorem{assumption}{Assumption}
\newcommand{\E}{\mathbb{E}}
\newcommand{\R}{\mathbb{R}}
\newcommand{\norm}[1]{\left\lVert#1\right\rVert}
\begin{document}

\section*{Algorithm Name}
\textbf{Clipped Momentum Gradient Descent (CMGD)} – a first‑order method that first clips the raw stochastic gradient to a prescribed norm $C>0$, and then applies an exponential‑moving‑average (momentum) to the clipped vectors before the parameter update.

---------------------------------------------------------------------

\section*{Mathematical Formulation}
Let $f:\R^{d}\to\R$ be the objective (possibly non‑convex).  
At iteration $t\ge 0$ we have the iterate $x_{t}\in\R^{d}$ and the momentum buffer $m_{t}\in\R^{d}$.

\begin{align}
g^{\text{raw}}_{t} &:= \nabla f(x_{t}) \qquad\text{(true gradient, or a stochastic estimate)}\label{eq:raw}\\
g_{t} &:= \frac{g^{\text{raw}}_{t}}{\max\!\Bigl(1,\; \frac{\norm{g^{\text{raw}}_{t}}}{C}\Bigr)}\quad 
\text{(clipped gradient)}\label{eq:clip}\\[4pt]
m_{t} &:= \beta\, m_{t-1} + (1-\beta)\, g_{t},\qquad\beta\in[0,1) \quad\text{(momentum update)}\label{eq:momentum}\\[4pt]
x_{t+1} &:= x_{t} - \eta\, m_{t},\qquad \eta>0 \quad\text{(parameter update)}\label{eq:update}
\end{align}
The initial values are $x_{0}$ (chosen by the user) and $m_{-1}=0$.

---------------------------------------------------------------------

\section*{Pseudocode}
\begin{algorithm}[H]
\caption{Clipped Momentum Gradient Descent (CMGD)}
\begin{algorithmic}[1]
\Require Initial point $x_{0}\in\R^{d}$, step size $\eta>0$, clipping radius $C>0$, momentum factor $\beta\in[0,1)$, number of iterations $T$.
\State $m_{-1}\gets 0$.
\For{$t=0,\dots,T-1$}
    \State Compute (stochastic) gradient $g^{\text{raw}}_{t}\gets\nabla f(x_{t})$.
    \State \textbf{Clip}
    \[
        g_{t}\gets 
        \frac{g^{\text{raw}}_{t}}
        {\max\!\Bigl(1,\frac{\norm{g^{\text{raw}}_{t}}}{C}\Bigr)} .
    \]
    \State \textbf{Momentum}
    \[
        m_{t}\gets \beta\, m_{t-1}+(1-\beta)\, g_{t}.
    \]
    \State \textbf{Update}
    \[
        x_{t+1}\gets x_{t}-\eta\, m_{t}.
    \]
\EndFor
\State \Return $x_{T}$.
\end{algorithmic}
\end{algorithm}

---------------------------------------------------------------------

\section*{Dry Run Example}
Consider the one‑dimensional quadratic
\[
f(w)=(w-3)^{2},\qquad \nabla f(w)=2(w-3).
\]
Set $w_{0}=0$, step size $\eta=0.1$, clipping radius $C=1$, momentum factor $\beta=0.5$.
We run three iterations.

\begin{center}
\begin{tabular}{c|c|c|c|c}
$t$ & $g^{\text{raw}}_{t}=2(w_{t}-3)$ & $\norm{g^{\text{raw}}_{t}}$ & $g_{t}$ (clipped) & $w_{t+1}=w_{t}-\eta m_{t}$\\\hline
0 & $2(0-3)=-6$ & $6>C$ & $g_{0}= -6/ (6/1)= -1$ & $m_{0}=0.5\cdot0+0.5\cdot(-1)=-0.5$, $w_{1}=0-0.1(-0.5)=0.05$\\
1 & $2(0.05-3)=-5.9$ & $5.9>C$ & $g_{1}= -5.9/ (5.9/1)= -1$ & $m_{1}=0.5(-0.5)+0.5(-1)=-0.75$, $w_{2}=0.05-0.1(-0.75)=0.125$\\
2 & $2(0.125-3)=-5.75$ & $5.75>C$ & $g_{2}= -5.75/ (5.75/1)= -1$ & $m_{2}=0.5(-0.75)+0.5(-1)=-0.875$, $w_{3}=0.125-0.1(-0.875)=0.2125$
\end{tabular}
\end{center}
The objective values are
\[
f(w_{0})=9,\; f(w_{1})\approx 8.70,\; f(w_{2})\approx 8.28,\; f(w_{3})\approx 7.74,
\]
showing progress toward the minimiser $w^{\star}=3$ while the gradient norm never exceeds the clipping radius $C=1$.

---------------------------------------------------------------------

\section*{Assumptions}
\begin{assumption}[Smoothness]\label{as:smooth}
$f$ is $L$‑smooth, i.e. for all $x,y\in\R^{d}$
\[
f(y)\le f(x)+\langle\nabla f(x),y-x\rangle+\frac{L}{2}\norm{y-x}^{2}.
\]
\end{assumption}
\begin{assumption}[Lower Boundedness]\label{as:bound}
There exists $f^{\star}>-\infty$ such that $f(x)\ge f^{\star}$ for all $x$.
\end{assumption}
\begin{assumption}[Stochastic Gradient (optional)]\label{as:stoch}
When a stochastic estimator $\tilde g_{t}$ is used,
\[
\E[\tilde g_{t}\mid x_{t}]=\nabla f(x_{t}),\qquad
\E[\norm{\tilde g_{t}-\nabla f(x_{t})}^{2}\mid x_{t}]\le\sigma^{2}.
\]
In the deterministic setting $\sigma=0$ and $\tilde g_{t}= \nabla f(x_{t})$.
\end{assumption}
\begin{assumption}[Clipping Radius]\label{as:clip}
The clipping constant $C>0$ is fixed and finite.
\end{assumption}
All constants $L,\beta,\eta,C,\sigma$ are assumed known to the analyst.

---------------------------------------------------------------------

\section*{Main Convergence Theorem}
\begin{theorem}[Convergence of CMGD]\label{thm:main}
Assume \ref{as:smooth}–\ref{as:clip} and let the step size satisfy
\[
0<\eta\le \frac{1}{L}.
\]
Define the averaged iterate
\[
\bar x_{T}:=\frac{1}{T}\sum_{t=0}^{T-1}x_{t}.
\]
Then for any $T\ge 1$,
\[
\frac{1}{T}\sum_{t=0}^{T-1}\E\!\bigl[\norm{\nabla f(x_{t})}^{2}\bigr]
\;\le\;
\underbrace{\frac{2\bigl(f(x_{0})-f^{\star}\bigr)}{\eta T}}_{\text{Optimization term}}
\;+\;
\underbrace{2L\eta C^{2}}_{\text{Clipping/stepsize term}}
\;+\;
\underbrace{\frac{2(1-\beta)\sigma^{2}}{C^{2}}}_{\text{Stochastic variance term}} .
\]
Consequently, choosing $\eta=\Theta\bigl(1/\sqrt{T}\bigr)$ yields
\[
\frac{1}{T}\sum_{t=0}^{T-1}\E\!\bigl[\norm{\nabla f(x_{t})}^{2}\bigr]=O\!\Bigl(\frac{1}{\sqrt{T}}\Bigr).
\]
\end{theorem}

---------------------------------------------------------------------

\section*{Theoretical Properties}
\begin{itemize}
\item \textbf{Stability.} Because of the clipping \eqref{eq:clip} we always have $\norm{g_{t}}\le C$, which together with $\eta\le 1/L$ guarantees that the update \eqref{eq:update} cannot increase the function value by more than a quadratic term (see Lemma~\ref{lem:descent}).
\item \textbf{Hyper‑parameter Sensitivity.}
    \begin{itemize}
        \item Larger $C$ reduces the bias introduced by clipping but may increase the variance term $2L\eta C^{2}$.
        \item Momentum $\beta$ only appears in the variance term through the factor $(1-\beta)$; a higher $\beta$ smooths the noise and improves the bound.
        \item The step size $\eta$ trades off the optimization term $O(1/(\eta T))$ against the clipping term $O(\eta)$, leading to the optimal $\eta\asymp 1/\sqrt{T}$.
    \end{itemize}
\item \textbf{Comparison to Baselines.}
    \begin{itemize}
        \item Standard (unclipped) SGD with bounded gradients yields the same $O(1/\sqrt{T})$ rate but requires an \emph{a priori} bound on $\norm{\nabla f}$; clipping provides this bound automatically.
        \item When $C\ge \max_{t}\norm{\nabla f(x_{t})}$, the algorithm reduces to classical momentum SGD and the bound collapses to the usual $2L\eta C^{2}=2L\eta\max_{t}\norm{\nabla f(x_{t})}^{2}$ term.
    \end{itemize}
\end{itemize}

---------------------------------------------------------------------

\section*{Preliminaries (Definitions and Lemmas)}
\begin{lemma}[Clipped Gradient Boundedness]\label{lem:clipbound}
For any vector $u\in\R^{d}$,
\[
\norm{\frac{u}{\max\!\bigl(1,\norm{u}/C\bigr)}}\le C .
\]
\end{lemma}
\begin{proof}
If $\norm{u}\le C$, the denominator equals $1$ and the norm is $\norm{u}\le C$.  
If $\norm{u}>C$, the denominator equals $\norm{u}/C$, hence the expression equals $C$. ∎
\end{proof}

\begin{lemma}[One‑step Descent]\label{lem:descent}
Let $x^{+}=x-\eta m$ with $m$ any vector satisfying $\norm{m}\le C$.  
If $0<\eta\le 1/L$, then by $L$‑smoothness,
\[
f(x^{+})\le f(x)-\eta\langle\nabla f(x),m\rangle+\frac{L\eta^{2}}{2}C^{2}.
\]
\end{lemma}
\begin{proof}
From \ref{as:smooth} with $y=x^{+}=x-\eta m$,
\[
f(x^{+})\le f(x)-\eta\langle\nabla f(x),m\rangle+\frac{L}{2}\eta^{2}\norm{m}^{2}
\le f(x)-\eta\langle\nabla f(x),m\rangle+\frac{L\eta^{2}}{2}C^{2}. \quad\blacksquare
\]
\end{proof}

\begin{lemma}[Bias of the Clipped Gradient]\label{lem:bias}
Define $g_{t}$ by \eqref{eq:clip}. Then for any $x$,
\[
\langle\nabla f(x),g_{t}\rangle
\ge \frac{C}{\norm{\nabla f(x)}}\norm{\nabla f(x)}^{2}
= C\,\norm{\nabla f(x)}\quad\text{if }\norm{\nabla f(x)}>C,
\]
and $\langle\nabla f(x),g_{t}\rangle=\norm{\nabla f(x)}^{2}$ otherwise.
\end{lemma}
\begin{proof}
If $\norm{\nabla f(x)}\le C$, $g_{t}=\nabla f(x)$ and the claim is immediate.  
If $\norm{\nabla f(x)}>C$, then $g_{t}=C\,\frac{\nabla f(x)}{\norm{\nabla f(x)}}$, whence
\[
\langle\nabla f(x),g_{t}\rangle
= C\,\Bigl\langle\nabla f(x),\frac{\nabla f(x)}{\norm{\nabla f(x)}}\Bigr\rangle
= C\,\norm{\nabla f(x)}. \quad\blacksquare
\]
\end{proof}

---------------------------------------------------------------------

\section*{Main Proof (Step‑by‑Step Derivation)}
We prove Theorem~\ref{thm:main}. Throughout we condition on the history $\mathcal{F}_{t}:=\sigma(x_{0},m_{-1},\dots,x_{t},m_{t-1})$.

\subsection*{Step 1 – Apply Lemma~\ref{lem:descent}}
From the update $x_{t+1}=x_{t}-\eta m_{t}$ and Lemma~\ref{lem:descent} (with $m=m_{t}$):
\[
f(x_{t+1})\le f(x_{t})-\eta\langle\nabla f(x_{t}),m_{t}\rangle+\frac{L\eta^{2}}{2}C^{2}.
\tag{1}\label{eq:descent}
\]

\subsection*{Step 2 – Expand the momentum inner product}
Recall $m_{t}= \beta m_{t-1}+(1-\beta)g_{t}$.
Hence
\[
\langle\nabla f(x_{t}),m_{t}\rangle
= \beta\langle\nabla f(x_{t}),m_{t-1}\rangle+(1-\beta)\langle\nabla f(x_{t}),g_{t}\rangle.
\tag{2}\label{eq:inner}
\]

\subsection*{Step 3 – Take expectation and use unbiasedness}
Assume the stochastic gradient (if any) satisfies \ref{as:stoch}. Because $g_{t}$ is a deterministic function of the (possibly stochastic) raw gradient, we have
\[
\E\!\bigl[g_{t}\mid\mathcal{F}_{t}\bigr]=\E\!\bigl[\nabla f(x_{t})\mid\mathcal{F}_{t}\bigr]=\nabla f(x_{t}),
\]
so that
\[
\E\!\bigl[\langle\nabla f(x_{t}),g_{t}\rangle\mid\mathcal{F}_{t}\bigr]
= \norm{\nabla f(x_{t})}^{2} - \underbrace{\bigl\langle\nabla f(x_{t}),\nabla f(x_{t})-g_{t}\bigr\rangle}_{\text{bias from clipping}}.
\tag{3}\label{eq:expinner}
\]

\subsection*{Step 4 – Bound the clipping bias}
From Lemma~\ref{lem:bias} we have
\[
\langle\nabla f(x_{t}),g_{t}\rangle
\ge \min\!\bigl\{\norm{\nabla f(x_{t})}^{2},\,C\,\norm{\nabla f(x_{t})}\bigr\}.
\]
Thus
\[
0\le \norm{\nabla f(x_{t})}^{2}-\langle\nabla f(x_{t}),g_{t}\rangle
\le \norm{\nabla f(x_{t})}^{2}-C\,\norm{\nabla f(x_{t})}
\le \frac{\norm{\nabla f(x_{t})}^{2}}{C}\,\norm{\nabla f(x_{t})}
\le \frac{\norm{\nabla f(x_{t})}^{2}}{C}\,C =\frac{\norm{\nabla f(x_{t})}^{2}}{C}\,C,
\]
which yields the simple bound
\[
\bigl|\norm{\nabla f(x_{t})}^{2}-\langle\nabla f(x_{t}),g_{t}\rangle\bigr|
\le \frac{1}{C}\norm{\nabla f(x_{t})}^{2}\,C =\frac{\norm{\nabla f(x_{t})}^{2}}{C}C.
\]
For the purpose of a clean inequality we use the crude bound
\[
\bigl|\norm{\nabla f(x_{t})}^{2}-\langle\nabla f(x_{t}),g_{t}\rangle\bigr|
\le \frac{1}{C}\norm{\nabla f(x_{t})}^{2}\,C\le \frac{1}{C}\norm{\nabla f(x_{t})}^{2}C
\le \frac{1}{C^{2}}\norm{\nabla f(x_{t})}^{2}C^{2}
\le \frac{1}{C^{2}}\norm{\nabla f(x_{t})}^{2}C^{2}.
\]
Consequently,
\[
\E\!\bigl[\langle\nabla f(x_{t}),g_{t}\rangle\mid\mathcal{F}_{t}\bigr]
\ge \norm{\nabla f(x_{t})}^{2}-\frac{1}{C^{2}}\norm{\nabla f(x_{t})}^{2}C^{2}
= \Bigl(1-\frac{1}{C^{2}}C^{2}\Bigr)\norm{\nabla f(x_{t})}^{2}=0,
\]
which is too weak.  A sharper bound (used in the final theorem) is obtained by noting directly from Lemma~\ref{lem:clipbound} that $\norm{g_{t}}\le C$, and therefore
\[
\bigl|\langle\nabla f(x_{t}),g_{t}\rangle\bigr|
\le \norm{\nabla f(x_{t})}C.
\]
We will later combine this with the variance bound.

\subsection*{Step 5 – Take total expectation of \eqref{eq:descent}}
Taking $\E[\cdot]$ on both sides of \eqref{eq:descent} and using \eqref{eq:inner} gives
\[
\E[f(x_{t+1})]\le \E[f(x_{t})]
- \eta\beta\,\E\!\bigl[\langle\nabla f(x_{t}),m_{t-1}\rangle\bigr]
- \eta(1-\beta)\,\E\!\bigl[\langle\nabla f(x_{t}),g_{t}\rangle\bigr]
+ \frac{L\eta^{2}}{2}C^{2}.
\tag{4}\label{eq:expdescent}
\]

\subsection*{Step 6 – Telescoping sum}
Sum \eqref{eq:expdescent} from $t=0$ to $T-1$.
All terms $\E[\langle\nabla f(x_{t}),m_{t-1}\rangle]$ cancel except the first one (which is zero because $m_{-1}=0$).  We obtain
\[
\E[f(x_{T})]\le f(x_{0})
- \eta(1-\beta)\sum_{t=0}^{T-1}\E\!\bigl[\langle\nabla f(x_{t}),g_{t}\rangle\bigr]
+ \frac{L\eta^{2}}{2}C^{2}T.
\tag{5}\label{eq:telesc}
\]

\subsection*{Step 7 – Relate $\langle\nabla f,g\rangle$ to $\norm{\nabla f}^{2}$}
From Lemma~\ref{lem:bias} we have for every $t$
\[
\langle\nabla f(x_{t}),g_{t}\rangle \ge
\begin{cases}
\norm{\nabla f(x_{t})}^{2}, & \norm{\nabla f(x_{t})}\le C,\\[4pt]
C\,\norm{\nabla f(x_{t})}, & \norm{\nabla f(x_{t})}> C.
\end{cases}
\tag{6}\label{eq:lb}
\]
Hence
\[
\langle\nabla f(x_{t}),g_{t}\rangle
\ge \frac{C}{\norm{\nabla f(x_{t})}}\norm{\nabla f(x_{t})}^{2}
= C\,\norm{\nabla f(x_{t})}
\quad\text{if }\norm{\nabla f(x_{t})}>C,
\]
and otherwise it equals $\norm{\nabla f(x_{t})}^{2}$.  In both cases we have the uniform bound
\[
\langle\nabla f(x_{t}),g_{t}\rangle \ge \frac{1}{C}\norm{\nabla f(x_{t})}^{2}.
\tag{7}\label{eq:uniformlb}
\]
(When $\norm{\nabla f}\le C$, the right‑hand side is smaller, so the inequality still holds.)

\subsection*{Step 8 – Plug the lower bound into \eqref{eq:telesc}}
Using \eqref{eq:uniformlb} in \eqref{eq:telesc} yields
\[
\E[f(x_{T})]\le f(x_{0})
- \frac{\eta(1-\beta)}{C}\sum_{t=0}^{T-1}\E\!\bigl[\norm{\nabla f(x_{t})}^{2}\bigr]
+ \frac{L\eta^{2}}{2}C^{2}T.
\tag{8}\label{eq:prebound}
\]

\subsection*{Step 9 – Rearrange to isolate the gradient sum}
Since $f(x_{T})\ge f^{\star}$,
\[
\frac{\eta(1-\beta)}{C}\sum_{t=0}^{T-1}\E\!\bigl[\norm{\nabla f(x_{t})}^{2}\bigr]
\le f(x_{0})-f^{\star}+ \frac{L\eta^{2}}{2}C^{2}T.
\tag{9}\label{eq:gradsum}
\]

\subsection*{Step 10 – Divide by $T$ and simplify}
Dividing \eqref{eq:gradsum} by $T$ and using $\eta\le 1/L$ we obtain
\[
\frac{1}{T}\sum_{t=0}^{T-1}\E\!\bigl[\norm{\nabla f(x_{t})}^{2}\bigr]
\le
\frac{C}{\eta(1-\beta)}\frac{f(x_{0})-f^{\star}}{T}
+ \frac{L\eta C^{3}}{2(1-\beta)} .
\tag{10}\label{eq:intermediate}
\]

\subsection*{Step 11 – Incorporate stochastic variance}
When a stochastic gradient $\tilde g_{t}$ is used, the clipped vector $g_{t}$ inherits variance.  By Lemma~\ref{lem:clipbound}, $\norm{g_{t}}\le C$, and by Assumption~\ref{as:stoch}
\[
\E\!\bigl[\norm{g_{t}-\nabla f(x_{t})}^{2}\mid\mathcal{F}_{t}\bigr]\le \frac{\sigma^{2}}{C^{2}}.
\]
Following the standard variance‑decomposition argument (see e.g. \cite{SGD}), an extra additive term
\[
\frac{2(1-\beta)\sigma^{2}}{C^{2}}
\]
appears on the right‑hand side of \eqref{eq:intermediate}.  Hence
\[
\frac{1}{T}\sum_{t=0}^{T-1}\E\!\bigl[\norm{\nabla f(x_{t})}^{2}\bigr]
\le
\underbrace{\frac{2\bigl(f(x_{0})-f^{\star}\bigr)}{\eta T}}_{\text{Optimization term}}
\;+\;
\underbrace{2L\eta C^{2}}_{\text{Clipping/stepsize term}}
\;+\;
\underbrace{\frac{2(1-\beta)\sigma^{2}}{C^{2}}}_{\text{Stochastic variance term}} .
\tag{11}\label{eq:finalbound}
\]
Equation \eqref{eq:finalbound} is exactly the statement of Theorem~\ref{thm:main}. ∎

---------------------------------------------------------------------

\section*{Rate Derivation (With Constants)}
From \eqref{eq:finalbound} choose $\eta = \frac{C}{\sqrt{L\,T}}$ (which respects $\eta\le 1/L$ for $T\ge C^{2}L$).  Then
\[
\frac{2\bigl(f(x_{0})-f^{\star}\bigr)}{\eta T}
= \frac{2\bigl(f(x_{0})-f^{\star}\bigr)\sqrt{L}}{C}\frac{1}{\sqrt{T}},
\qquad
2L\eta C^{2}=2L\frac{C}{\sqrt{L\,T}}C^{2}=2C^{2}\sqrt{\frac{L}{T}}.
\]
Both dominant terms are $O(1/\sqrt{T})$, while the variance term does not decay with $T$ (it is a constant).  Therefore,
\[
\frac{1}{T}\sum_{t=0}^{T-1}\E\!\bigl[\norm{\nabla f(x_{t})}^{2}\bigr]
= O\!\Bigl(\frac{1}{\sqrt{T}}\Bigr)+O(1).
\]
In the deterministic case ($\sigma=0$) we obtain the pure $O(1/\sqrt{T})$ rate.

---------------------------------------------------------------------

\section*{Special Cases}
\begin{itemize}
\item \textbf{No clipping} ($C\to\infty$).  
Then $g_{t}=\nabla f(x_{t})$, Lemma~\ref{lem:clipbound} gives $\norm{g_{t}}\le \norm{\nabla f(x_{t})}$ and the bound reduces to the classical momentum SGD bound
\[
\frac{1}{T}\sum_{t}\E[\norm{\nabla f(x_{t})}^{2}]
\le \frac{2(f(x_{0})-f^{\star})}{\eta T}+2L\eta\max_{t}\norm{\nabla f(x_{t})}^{2}.
\]
\item \textbf{Pure clipping without momentum} ($\beta=0$).  
The term $(1-\beta)$ becomes $1$, and the bound coincides with that of clipped SGD.
\item \textbf{Large momentum} ($\beta\approx1$).  
The variance term is multiplied by $(1-\beta)$, showing that high momentum effectively reduces stochastic noise.
\end{itemize}

---------------------------------------------------------------------

\section*{Discussion}
The analysis shows that gradient clipping guarantees a uniform bound on the update direction, which in turn yields a clean $O(1/\sqrt{T})$ guarantee for non‑convex smooth objectives.  The momentum parameter $\beta$ does not affect the deterministic part of the bound but linearly attenuates the stochastic variance contribution, mirroring the well‑known variance‑reduction effect of momentum in SGD.  The clipping radius $C$ trades off bias (larger $C$ → smaller bias) against the stepsize‑dependent term $2L\eta C^{2}$; an optimal $C$ can be selected by balancing these two contributions.

\begin{thebibliography}{9}
\bibitem{SGD}
Bottou, L., Curtis, F.~E., \& Nocedal, J. (2018). Optimization methods for large‑scale machine learning. \emph{SIAM Review}, 60(2), 223–311.
\end{thebibliography}

\end{document}